% ============================================
% 第2章：为什么孩子会蹲在沙坑里搭三个小时的桥
% 梁桥 · 游戏本质
% ============================================

\chapter{为什么孩子会蹲在沙坑里搭三个小时的桥}

\vspace{0.5cm}

\begin{center}
{\small\color{muted} 一个学完了材料力学、结构力学、理论力学的大学生，}
{\small\color{muted} 当他第一次打开 Poly Bridge 搭桥的时候，}
{\small\color{muted} 他的行为和一个六岁小孩完全一样。}
\end{center}

\vspace{1cm}

\section{洗浴中心里的桥梁工程师}

2026年5月22日，我在一场面向呼和浩特铁路局教师的线上分享会上，讲了一段话，后来被转写成了文字：

\begin{shiyu-inline}
``学生在洗浴的时候拿个平板，然后一直在玩这个游戏，一直在建不同的桥梁，玩到这个洗浴都关门了。那可能九十点钟还在这个沉浸的这个过程。''

\hfill ——2026年5月22日，游戏化教学分享，录音转写
\end{shiyu-inline}

一个大学生，在洗浴中心，拿着平板，搭桥搭到洗浴关门。

这不是一个段子。这是真实发生的事情。我的学生，在大连的某家洗浴中心，泡完澡之后，没有刷短视频，没有打王者荣耀，而是在用 Poly Bridge 搭桥。他在搭一座斜拉桥。他调了三次缆索的张力，改了两次桥塔的高度，最后一辆小车稳稳地开了过去。他抬头一看，洗浴中心要关门了。

我后来问过他：``你当时在想什么？''

他说：``没想什么。就是塌了，想让它站住。''

\section{``没有人说我用AI去糊弄帮我打游戏''}

在同一场分享里，我还说了另一句话：

\begin{shiyu-inline}
``没有人说我用AI去糊弄帮我打游戏。如果先去下棋或者打牌，没有人说我去用AI来帮我下棋或者打牌。因为我们本身就喜欢这种活动。那我们能不能把这个打牌呀、打游戏呀这种状态，让他们在课堂上也有这种状态，然后就是沉浸在这里面，不会去想我怎么去糊弄他，而是真正的喜欢他。''

\hfill ——2026年5月22日，同一场分享
\end{shiyu-inline}

这句话里有一个特别重要的对比：\textbf{学生会用AI糊弄一份作业，但没有人会用AI糊弄一场游戏。}

为什么？因为作业是别人让你做的事。游戏是你自己选择做的事。作业的目标是``完成''。游戏的目标是``赢''——或者说，是``经历从不会到会的过程，并且在这个过程中感受到自己正在变强''。

这不是我的理论。这是任何一个在洗浴中心搭桥搭到关门的大学生，用他的身体告诉我的。

\section{``很热闹，但是学生没学到什么''}

但游戏化不是没有争议的。

2026年5月31日，我做了另一场游戏化教学分享。开场前，我用一个``六顶思考帽''的活动，收集了老师们对游戏化的看法。后来这些话被转写成了文字：

\begin{shiyu-inline}
``游戏化有哪些价值、好处和收益呢？提高学生积极性啊，形式上比单纯讲授新颖一点，体验性强一点，兴趣是最好的老师，吸引学生的兴趣，提升课堂质量，让这个孩子做课堂的主人。''

``消极的部分——课堂纪律难以把控。学习重点可能会偏移。很热闹，但是学生没学到什么。玩到位了，是不是学到位了？还有就是，领导可能会不喜欢，影响考核。学生太活跃，领导不喜欢，影响我考核，降低我薪资。''

\hfill ——2026年5月31日，游戏化教学设计分享，录音转写
\end{shiyu-inline}

这些话，我每次分享都会听到。它们分成两派：

\textbf{积极派说：}游戏化能提高兴趣、增加参与、让学生做课堂的主人。

\textbf{消极派说：}游戏化会让课堂失控、目标偏移、热闹但没学到东西、领导不喜欢。

两派说的都对。但两派都漏掉了一个问题：\textbf{游戏为什么好玩？}

\section{游戏的核心：即时反馈}

2026年5月22日，我在讲 Poly Bridge 的时候，说了一段话：

\begin{shiyu-inline}
``其实我们非常关注的就是游戏它本身有什么特征，跟我们传统的教学有什么本质的区别。其实我感觉最重要的就是\textbf{及时反馈}。就是我做出来一个操作，马上我就能获得这个东西行还是不行，好还是不好，这样的一个回应。然后我跟这个回应，再去做这个调试。''

\hfill ——2026年5月22日，游戏化教学分享
\end{shiyu-inline}

这句话是我对游戏本质的理解：\textbf{即时反馈。}

你动一下，桥马上有反应。你加一根杆，桥的颜色变了。你放一辆小车，桥塌了——不是考试之后才知道，是\textbf{当场就知道}。你改了三个地方，又放一辆小车，桥站住了。你心里``啊''了一声。

这个回路——\textbf{动作→反馈→调整→再动作}——在游戏里，几秒钟就能完成一轮。在传统的课堂里，这个回路可能需要几周。你听了一节课，做了一次作业，交上去，等老师批改，拿回来，看到分数——但你已经忘了你当时是怎么想的了。

\subsection{失败的权利}

在同一场分享里，我还说了一句后来被反复引用的话：

\begin{shiyu-inline}
``学生在学习过程中从来没有失败的权利，他也就没有成功的那个喜悦或者荣耀。''

\hfill ——2026年5月22日，同一场分享
\end{shiyu-inline}

这句话是我在展示一段视频时说的。视频里，一座桥塌了。然后学生改了三个地方，桥站住了。整个过程不到两分钟。

在传统的课堂里，失败是\textbf{被惩罚的}。你答错了，扣分。你作业做错了，扣分。你考试没考好，扣分。但在游戏里，失败是\textbf{免费的}。你塌了一百次桥，没有人扣你的分。你只需要再试一次。每一次失败都告诉你一件事：\textbf{这里不对，换一种方式。}

\begin{insight}
\label{insight:chap2}

\textbf{游戏不是逃避。游戏是主动选择困难，并在安全的环境中反复失败，直到成功。}

孩子在沙坑里搭三个小时的桥，不是因为好玩——是因为他正在经历一个他主动选择的困难，并且他正在克服它。每一次沙堡塌了，他都知道为什么塌。每一次沙堡站住了，他都感到一种发自心底的骄傲。

这种骄傲，考试得100分给不了。
\end{insight}

\section{``我们从来没有教过学生从零开始跟工程师一起去思考''}

2026年5月31日，我讲了另一段话：

\begin{shiyu-inline}
``我们传统的教学里面从来没有教过学生从零开始跟工程师一起去思考这个问题。之前我们就告诉说，应该是这么做的是吧？这个桁架一般这么高，五米高，二十米的桥应该是五米高的桁架。这个构件应该什么尺寸。这个状态是很可悲的——因为他在设计别人做的结构。他从来没有真正像一个工程师一样，说我们思考一下这个问题该怎么解决。它可能不一定只有一个解法。''

\hfill ——2026年5月31日，游戏化教学设计分享
\end{shiyu-inline}

这段话里有一个词，我自己说出来的时候都觉得有点重：\textbf{可悲}。

我们教学生``应该怎么做''。我们告诉他们桁架应该多高，梁应该多宽，配筋应该怎么配。我们从来没有让他们自己想过：\textbf{如果让你来设计，你会怎么做？}

但在游戏里，学生第一次被放在了\textbf{设计者}的位置上。不是``记住这个公式，考试会考''，而是``这里有一道沟，你有一座桥要搭，成本不能超过三万，车要能过，你自己想办法''。

这个区别，就是\textbf{学习}和\textbf{被灌输}之间的区别。

\section{``六岁小孩''和``大学生''}

2026年5月22日，我说了那句话：

\begin{shiyu-inline}
``学生在上这门课之前已经学过材料力学、结构力学、或者理论力学，但是他搭建这个桥梁的过程或者状态，跟他六岁的小孩搭建这个状态是完全一致的。我们不是说他这个多么低龄或多么幼稚。我们在说，其实我们讲授的很多的理论的这种课程，如果没有经过这种实操、试错、迭代、失败成功的这个过程，他可能很难入到你一个学生的这个直觉里面去。''

\hfill ——2026年5月22日，游戏化教学分享
\end{shiyu-inline}

这句话里有一个非常重要的词：\textbf{直觉}。

一个学完了材料力学、结构力学、理论力学的大学生，当他第一次搭桥的时候，他的行为和一个六岁小孩\textbf{完全一样}。他不是在调用公式，他是在用直觉——这里加一根杆试试，那里加高一点试试，塌了，换一种方式。

这不是退化。这是\textbf{学习的起点}。所有的理论，所有的公式，所有的规范，都必须经过这个起点——\textbf{动手做，失败，调整，再动手}——才能真正进入一个人的直觉。没有这个过程，你背再多的公式，那个公式也只是黑板上的一行字，不是你手指上的感觉。

\section{梁桥的隐喻之二}

梁桥是最简单的桥。但最简单的，最难做好。因为梁桥的跨度越大，自重越大——80\%的自重，20\%的荷载。

\textbf{游戏的本质，就是帮我们解决梁桥的困境。}

游戏不是``加糖''。游戏不是``让学习变轻松''。游戏是\textbf{重新设计学习回路}——把困惑→尝试→失败→调整→成功的循环，从几周压缩到几分钟。让学生在安全的环境中反复失败，直到他们自己发现那个公式的含义。

当一个学生蹲在沙坑里搭三个小时的桥，他不是在``玩''。他是在\textbf{学习}——用人类最古老、最自然、最有效的方式学习。

\begin{shiyu-note}
\label{note:chap2}

\textbf{游戏不是学习的敌人。游戏是学习最原始的形态。}

我们总是把``学习''和``严肃''捆绑在一起。好像不够严肃的东西，就不可能是学习。但所有三岁孩子都知道：学习是一件让人上瘾的事。你学走路，摔倒了，爬起来，再摔，再爬起来——你从来没有觉得``学走路好累，我不学了''。因为走路不是别人让你学的，是\textbf{你自己想学的}。

游戏化教学，不是把课堂变成游乐场。是把课堂重新设计成\textbf{一个可以安全失败、即时反馈、主动探索的地方}。就像沙坑。就像 Poly Bridge。就像任何一个孩子蹲在那里搭了三个小时的桥，然后抬头说：``站住了。''

\end{shiyu-note}

\vspace{1cm}

\begin{decision-point}
\label{decision:chap2}

\noindent\textbf{这一章，我们看到了游戏的本质：即时反馈、安全失败、主动选择困难。}

\vspace{0.3cm}

\noindent 现在，回顾你上一次批改的作业。你花了多少时间批改？学生拿到批改结果之后，距离他做作业的时候，已经过了多久？

\noindent 如果这个反馈回路太长，试着缩短它。下一节课，不要等到课后批改。在课堂上，让学生当场做，当场看结果，当场改。就一次。

\noindent 告诉我，发生了什么。
\end{decision-point}

\vspace{0.5cm}

\begin{flushright}
{\small\color{muted} 下一章：当学生带着ChatGPT的答案走进教室——旧系统是如何被冲击的。}
\end{flushright}